\subsection{Interpretability analysis}

To better understand and interpret the behavior of FFNNs, which are often regarded as ``black-box'' models due to their complex and non-linear structure, we employed two explainable AI (xAI) techniques~\cite{BARREDOARRIETA202082}: partial dependence plots (PDPs)~\cite{friedman2001pdps} and Shapley values~\cite{Merrick2020explaination}.
These methods allow us to quantify the influence of individual predictors on model outputs and to visualize their effects across the predictor space.
In particular, we used them to assess the robustness of BELLA to irrelevant input predictors, as well as to gain general insights into how the models leverage relevant features in the inference task.

Partial dependence plots (PDPs) illustrate the average effect of a single predictor \(x_S\) (or a subset of predictors) on the model’s output, marginalizing over the distribution of all other features \(x_C\). 
Formally, for a prediction function \(f(\mathbf{x})\), the partial dependence of the feature subset \(S\) is defined as
\[
    \mathrm{PD}_S(x_S) = \mathbb{E}_{\mathbf{x}_C} \big[ f(x_S, \mathbf{x}_C) \big],
\]
where \( C \) denotes the complement of \( S \), and the expectation is taken with respect to the marginal distribution of the remaining features, \( p(\mathbf{x}_C) \). 

PDPs are particularly useful because they provide an interpretable, model-agnostic visualization of how an input feature influences the prediction on average, even in the presence of complex nonlinear interactions captured by models such as FFNNs.
If the predictor is relevant, the PDP typically exhibits a structured relationship (e.g., monotonic trend, threshold effect, or nonlinear curve).
Conversely, if the predictor is irrelevant, the PDP should remain approximately flat, indicating that varying this feature does not systematically alter the model output.
This makes PDPs well-suited for diagnosing spurious dependencies and detecting overfitting when irrelevant predictors are included in the input space.
However, PDPs also have important limitations, as they implicitly assume independence between features when marginalizing, which may lead to misleading plots if predictors are correlated.

Shapley addresses these issues by offering a game-theoretic framework for feature attribution.
For a model \( f \) and input instance \( \mathbf{x} \), the Shapley value for feature \( i \) is:
\[
    \phi_i(f, \mathbf{x}) = \sum_{S \subseteq F \setminus \{i\}}
    \frac{|S|! \, (|F|-|S|-1)!}{|F|!}
    \Big[ f_{S \cup \{i\}}(\mathbf{x}_{S \cup \{i\}}) - f_S(\mathbf{x}_S) \Big],
\]
where \( F \) is the full feature set and \( S \) a subset excluding \( i \), and the term \( f_S(\mathbf{x}_S) \) is defined through a conditional expectation:
\[
    f_S(\mathbf{x}_S) = \mathrm{E}_{\mathbf{x}_{F \setminus S}}\!\left[f(\mathbf{x}_S, \mathbf{x}_{F \setminus S}) \mid \mathbf{x}_S \right],
\]
that is, the model output averaged over the distribution of the missing features given the known ones.

Shapley values provide local, prediction-level explanations, but in practice, feature importances are often summarized by averaging the absolute Shapley values across all instances in the input dataset:
\[
    \Phi_i(f) = \frac{1}{|X|} \sum_{j=1}^{|X|} |\phi_i(f, \mathbf{x}^{(j)})|,
\]
where \( |X| \) is the number of samples.
This provides a global quantification of each feature’s contribution, complementing the local and marginal insights offered by PDPs, and enables a more comprehensive assessment of how predictors influence the network outputs across different conditions.

In our study, each FFNN has multiple configurations due to the weights sampled during the MCMC process.
To compute PDPs and Shapley values, we first subsampled 100 MCMC samples per run, yielding 100 distinct FFNN models per tree.
For each input predictor, we computed the PDP values for each model, and then aggregated them into a single PDP per tree by taking their median.
Similarly, the SHAP framework \cite{lundberg2017unified} was used to compute Shapley feature importance scores for each of the 100 subsampled FFNN models; these were then aggregated by taking the median feature importance per model and finally normalized so that feature contributions per model sum to 1.